\documentclass[aspectratio=169,12pt]{beamer}

%----------------------------------------------------
% Packages
%----------------------------------------------------
\RequirePackage{Gabby} % optional, if available
\usepackage{xcolor}
\usepackage{graphicx}
\usepackage{amsmath,amssymb}
\usepackage{natbib}

%----------------------------------------------------
% Custom Command
%----------------------------------------------------
\newcommand{\bcol}[2][blue]{\textcolor{#1}{#2}}

%----------------------------------------------------
% Title Information (EDIT THIS)
%----------------------------------------------------
\title[Project Title (Short)]%
{Full Project Title Goes Here}

\author[Student Names]%
{{\scriptsize by}\\
	\bcol[green!40!black]{Student Name 1 \quad Student Name 2 \quad Student Name 3}\\
	\bigskip
	{\small Supervisor: Prof.\ Supervisor’s Name}}

\institute[KNUST-Maths]%
{Department of Mathematics\\
	Kwame Nkrumah University of Science and Technology}

\logo{\includegraphics[scale=0.1]{KNUSTlogo}}
\date{\scriptsize \today}

%====================================================
\begin{document}
	
	%----------------------------------------------------
	% Title Page
	%----------------------------------------------------
	\begin{frame}
		\titlepage
	\end{frame}
	
	%----------------------------------------------------
	% Outline
	%----------------------------------------------------
	\begin{frame}{Outline of Presentation}
		\tableofcontents
	\end{frame}
	
	%====================================================
	\section{Background of the Study  (Maximum 4 Slides)}
	
	\begin{frame}{Background of the Study}
		\begin{block}{Purpose of This Section}
			Introduce the general area of study and provide the context of the research.
		\end{block}
		
		\begin{itemize}
			\item Briefly describe the real-world or academic problem.
			\item Introduce key concepts needed to understand the project.
			\item Avoid technical details at this stage.
		\end{itemize}
		
		\begin{alertblock}{Presentation Tip}
			Use clear bullet points. Each point should be understandable without further explanation.
		\end{alertblock}
		
		\begin{exampleblock}{Good Practice}
			Where possible, use figures, diagrams, or illustrations instead of long text.
		\end{exampleblock}
	\end{frame}
	
	%====================================================
	\section{Problem Statement  (Maximum 2 Slides)}
	
	\begin{frame}{Problem Statement}
		\begin{columns}[T]
			\begin{column}{0.48\textwidth}
				\begin{block}{What Is the Problem?}
					\begin{itemize}
						\item Clearly state the specific problem being addressed.
						\item Explain what is missing or inadequate in existing studies.
					\end{itemize}
				\end{block}
			\end{column}
			
			\begin{column}{0.48\textwidth}
				\begin{block}{Why Is It a Problem?}
					\begin{itemize}
						\item Explain why the problem is important.
						\item Relate it to real-world applications or theory.
					\end{itemize}
				\end{block}
			\end{column}
		\end{columns}
		
		\begin{alertblock}{Presentation Tip}
			Do not read directly from the slides. Explain the problem in your own words.
		\end{alertblock}
	\end{frame}
	
	%====================================================
	\section{Objectives  (Maximum 1 Slide)}
	
	\begin{frame}{Objectives of the Study}
		\begin{block}{General Objective}
			State the main aim of the project in one sentence.
		\end{block}
		
		\begin{block}{Specific Objectives}
			\begin{enumerate}
				\item Begin each objective with an action verb (e.g., develop, analyse, investigate).
				\item Objectives should be achievable and measurable.
				\item Avoid equations unless absolutely necessary.
			\end{enumerate}
		\end{block}
	\end{frame}
	
	%====================================================
	\section{Methodology  (Maximum 8 Slides)}
	
	\begin{frame}{Methodology}
		\begin{itemize}
			\item Describe the approach used to solve the problem. \pause
			\item Present models, algorithms, or analytical techniques.  \pause
			\item Explain assumptions clearly.
		\end{itemize}
		
		\begin{exampleblock}{Mathematics Projects}
			Clearly define variables, parameters, and governing equations.
		\end{exampleblock}
	\end{frame}
	
	\begin{frame}{Referencing}
		\begin{block}{Citation Style}
			Referencing in Beamer follows the same style as the main project report.
		\end{block}
		
		\begin{enumerate}
			\item In-text citation: \cite{Fosu2020}
			\item Parenthetical citation: \citep{Appati2015}
			\item Multiple sources: \citep{Fosu2021,Appati2015}
		\end{enumerate}
		
\vspace{1cm}

		Follow this link for more on referencing with Jabref
\begin{center}
\href{https://youtu.be/zzlXj93WgJQ}{\bcol{https://youtu.be/zzlXj93WgJQ}}

\end{center}
		
	\end{frame}
	
	%====================================================
	\section{Analysis and Results  (Maximum 8 Slides)}
	
	\begin{frame}{Analysis and Results}
		\begin{itemize}[<+->]
			\item Present key results only (avoid overcrowding slides).
			\item Use tables and figures where possible.
			\item Number equations for easy referencing.
		\end{itemize}
		
		\begin{alertblock}{Important}
			Always interpret results - do not just display them.
		\end{alertblock}
	\end{frame}
	
	
	\begin{frame}[fragile]{Some Tips} 
		\begin{enumerate} 
			\item Working with Tables, Figures and Equations is the same for beamer. 
			\item Bullet your points: Each point should be powerful enough to explain itself. 
			\item If possible, use pictures to introduce the background. 
			\item Capture literature as part of the background. 
			\item Ensure proper spacing by using 
			\begin{verbatim} 
				\medskip, \bikskip, \vspace{2cm}, \qquad, \quad 
			\end{verbatim} 
			\item Colours could be used to highlight some equations 
			\begin{equation}\label{key} \bcol[red]{2+ \sin x = 2000} \qquad \bcol[green]{2+ \sin x = 2000}\qquad \bcol{2+ \sin x = 2000} 
		\end{equation} 
	\end{enumerate} 
\end{frame}
	
	%====================================================
	\section{Conclusion and Recommendations  (Maximum 3 Slides)}
	
	\begin{frame}{Conclusion and Recommendations}
		\begin{block}{Conclusion}
			Summarise the main findings and relate them to the objectives.
		\end{block}
		
		\begin{block}{Recommendations}
			Suggest practical applications, improvements, or directions for future work.
		\end{block}
	\end{frame}
	
	%====================================================
	\section{Selected References}
	
	\begin{frame}{Selected References}
		\small
		\bibliographystyle{apalike}
		\bibliography{myref}
	\end{frame}
	
	%====================================================
	\section*{}
	
	\begin{frame}
		\begin{center}
			{\Huge Thank You}
		\end{center}
	\end{frame}
	
	%====================================================
\end{document}
